\section{Introduction}\label{sec:intro}

\subsection{The question}

Let $f(x)=x+1$ and $g(x)=-1/x$. Start at $0$ and apply $f$ and $g$ in any order, never applying $g$
to $0$. Every rational is reached, and the \emph{rank} of a rational is the least number of moves
needed to reach it. Kimberling sorted the rationals by rank in OEIS A226247 \cite{oeisA226247}, and
Kagey asked whether the number $a(n)$ of rationals of rank $n$ satisfies $a(n)=a(n-1)+a(n-3)$
\cite{kagey-mse,kagey137}. In earlier work \cite{P137} we proved this for $n\ge4$, together with a
formula for the rank and a description of the tree. Kimberling and Moses had stated the same
recurrence for the version of the tree with root $1$ \cite[Example~4.2]{kimberling-moses}.

This paper asks where the $3$ comes from. The answer is short. The map
\[
g(f(x))=-\frac{1}{x+1}
\]
has order $3$, since $x\mapsto-1/(x+1)\mapsto-(x+1)/x\mapsto x$. To test this reading we change $g$
so that $g\circ f$ has a different order and see what happens to the recurrence. Keep $f(x)=x+1$ and
the root $0$, and for an integer $m\ge1$ put
\[
g_m(x)=-\frac{1}{mx}.
\]
Problem 137 is the case $m=1$. The composite $g_m(f(x))=-1/(m(x+1))$ is a M\"obius map with
trace $m$ and determinant $m$. It has order $3$, $4$ and $6$ for $m=1$, $2$ and $3$, and infinite
order for $m\ge4$ (Lemma~\ref{lem:order}). We write $q$ for this order. Table~\ref{tab:family}
summarizes what we prove about each case.

\begin{table}[ht]
\centering\small
\renewcommand{\arraystretch}{1.15}
\begin{tabular}{@{}c c c l l@{}}
\toprule
$m$ & $q$ & points of the tree & generating function of $a(n)$ & recurrence\\
\midrule
$1$ & $3$ & all of $\Q$ & $\dfrac{1+t^2-t^3}{1-t-t^3}$ & $a(n)=a(n-1)+a(n-3)$, $n\ge4$\\[8pt]
$2$ & $4$ & all of $\Q$ & $\dfrac{1+t^2+t^4-t^5}{1-t-t^3-t^5}$ & $a(n)=a(n-1)+a(n-3)+a(n-5)$, $n\ge6$\\[8pt]
$3$ & $6$ & all of $\Q$ & $\dfrac{1+t^2+t^4+t^6+t^8-t^9}{1-t-t^3-t^5-t^7-t^9}$ &
  $a(n)=\sum_{j=0}^{4}a(n-2j-1)$, $n\ge10$\\[8pt]
$\ge4$ & $\infty$ & not all of $\Q$ & $\dfrac{1}{1-t-t^2}$ & $a(n)=a(n-1)+a(n-2)$, $n\ge2$\\
\bottomrule
\end{tabular}
\caption{The trees of $f(x)=x+1$ and $g_m(x)=-1/(mx)$ with root $0$. The row $m=1$ is
\cite{P137}. In the third row the sum is $a(n-1)+a(n-3)+a(n-5)+a(n-7)+a(n-9)$. For $m\le3$ each
recurrence fails at the first $n$ below the stated range.}
\label{tab:family}
\end{table}

So the $3$ in Problem 137 is the order of $g\circ f$, and the recurrence for order $q$ has the
$q-1$ lags $1,3,5,\dots,2q-3$. The family also shows why the more obvious generalization
$x\mapsto x+k$ with $x\mapsto-1/x$ is not interesting. Putting $x=ky$ turns it into the tree of
$y\mapsto y+1$ and $y\mapsto-1/(k^2y)$, which is the case $m=k^2\ge4$, and its counts are
Fibonacci numbers.

The cases $m=2$ and $m=3$ give new trees of rationals. They are the only other values of $m$ for
which the tree is all of $\Q$ (Theorem~D and Theorem~E). Their counts start
\begin{align*}
(m=2)&\qquad 1, 1, 2, 3, 5, 7, 11, 18, 28, 44, 69, 108, 170, 267, 419, 658, 1033, \dots,\\
(m=3)&\qquad 1, 1, 2, 3, 5, 8, 13, 21, 34, 54, 87, 141, 227, 366, 590, 951, 1533, \dots,
\end{align*}
and neither sequence is in the OEIS. Figure~\ref{fig:trees} shows the first ranks of the trees for
$m=1$ and $m=2$.

\begin{figure}[p]
\centering
\resizebox{0.86\textwidth}{!}{\begin{tikzpicture}[every node/.style={font=\scriptsize}]
\node[vroot] (n0_1) at (2.99,0.00) {$0$};
\node[vpos] (n1_1) at (2.99,-1.05) {$1$};
\node[vneg] (nm1_1) at (0.00,-2.10) {$-1$};
\node[vpos] (n2_1) at (5.98,-2.10) {$2$};
\node[vneg] (nm1_2) at (2.34,-3.15) {$-\tfrac{1}{2}$};
\node[vpos] (n3_1) at (9.61,-3.15) {$3$};
\node[vpos] (n1_2) at (2.34,-4.20) {$\tfrac{1}{2}$};
\node[vneg] (nm1_3) at (7.19,-4.20) {$-\tfrac{1}{3}$};
\node[vpos] (n4_1) at (12.03,-4.20) {$4$};
\node[vneg] (nm2_1) at (1.25,-5.25) {$-2$};
\node[vpos] (n3_2) at (3.44,-5.25) {$\tfrac{3}{2}$};
\node[vpos] (n2_3) at (7.19,-5.25) {$\tfrac{2}{3}$};
\node[vneg] (nm1_4) at (10.62,-5.25) {$-\tfrac{1}{4}$};
\node[vpos] (n5_1) at (13.44,-5.25) {$5$};
\node[vneg] (nm2_3) at (2.50,-6.30) {$-\tfrac{2}{3}$};
\node[vpos] (n5_2) at (4.38,-6.30) {$\tfrac{5}{2}$};
\node[vneg] (nm3_2) at (6.25,-6.30) {$-\tfrac{3}{2}$};
\node[vpos] (n5_3) at (8.12,-6.30) {$\tfrac{5}{3}$};
\node[vpos] (n3_4) at (10.62,-6.30) {$\tfrac{3}{4}$};
\node[vneg] (nm1_5) at (12.50,-6.30) {$-\tfrac{1}{5}$};
\node[vpos] (n6_1) at (14.38,-6.30) {$6$};
\node[vpos] (n1_3) at (2.50,-7.35) {$\tfrac{1}{3}$};
\node[vneg] (nm2_5) at (3.75,-7.35) {$-\tfrac{2}{5}$};
\node[vpos] (n7_2) at (5.00,-7.35) {$\tfrac{7}{2}$};
\node[vneg] (nm3_5) at (7.50,-7.35) {$-\tfrac{3}{5}$};
\node[vpos] (n8_3) at (8.75,-7.35) {$\tfrac{8}{3}$};
\node[vneg] (nm4_3) at (10.00,-7.35) {$-\tfrac{4}{3}$};
\node[vpos] (n7_4) at (11.25,-7.35) {$\tfrac{7}{4}$};
\node[vpos] (n4_5) at (12.50,-7.35) {$\tfrac{4}{5}$};
\node[vneg] (nm1_6) at (13.75,-7.35) {$-\tfrac{1}{6}$};
\node[vpos] (n7_1) at (15.00,-7.35) {$7$};
\node[rk] at (-0.6,0.63) {rank};
\node[rk] at (-0.6,0.00) {0};
\node[rk] at (-0.6,-1.05) {1};
\node[rk] at (-0.6,-2.10) {2};
\node[rk] at (-0.6,-3.15) {3};
\node[rk] at (-0.6,-4.20) {4};
\node[rk] at (-0.6,-5.25) {5};
\node[rk] at (-0.6,-6.30) {6};
\node[rk] at (-0.6,-7.35) {7};
\draw[ef] (n0_1) -- (n1_1);
\draw[eg] (n1_1) -- (nm1_1);
\draw[ef] (n1_1) -- (n2_1);
\draw[eg] (n2_1) -- (nm1_2);
\draw[ef] (n2_1) -- (n3_1);
\draw[ef] (nm1_2) -- (n1_2);
\draw[eg] (n3_1) -- (nm1_3);
\draw[ef] (n3_1) -- (n4_1);
\draw[eg] (n1_2) -- (nm2_1);
\draw[ef] (n1_2) -- (n3_2);
\draw[ef] (nm1_3) -- (n2_3);
\draw[eg] (n4_1) -- (nm1_4);
\draw[ef] (n4_1) -- (n5_1);
\draw[eg] (n3_2) -- (nm2_3);
\draw[ef] (n3_2) -- (n5_2);
\draw[eg] (n2_3) -- (nm3_2);
\draw[ef] (n2_3) -- (n5_3);
\draw[ef] (nm1_4) -- (n3_4);
\draw[eg] (n5_1) -- (nm1_5);
\draw[ef] (n5_1) -- (n6_1);
\draw[ef] (nm2_3) -- (n1_3);
\draw[eg] (n5_2) -- (nm2_5);
\draw[ef] (n5_2) -- (n7_2);
\draw[eg] (n5_3) -- (nm3_5);
\draw[ef] (n5_3) -- (n8_3);
\draw[eg] (n3_4) -- (nm4_3);
\draw[ef] (n3_4) -- (n7_4);
\draw[ef] (nm1_5) -- (n4_5);
\draw[eg] (n6_1) -- (nm1_6);
\draw[ef] (n6_1) -- (n7_1);
\draw[eold] (nm2_1) -- (nm1_1);
\draw[eold] (nm3_2) -- (nm1_2);
\draw[eold] (nm4_3) -- (nm1_3);
\end{tikzpicture}
}

\medskip
{\small (a) $m=1$, the tree of Problem 137.}

\bigskip
\resizebox{\textwidth}{!}{\begin{tikzpicture}[every node/.style={font=\scriptsize}]
\node[vroot] (n0_1) at (4.99,0.00) {$0$};
\node[vpos] (n1_1) at (4.99,-1.05) {$1$};
\node[vneg] (nm1_2) at (1.09,-2.10) {$-\tfrac{1}{2}$};
\node[vpos] (n2_1) at (8.88,-2.10) {$2$};
\node[vpos] (n1_2) at (1.09,-3.15) {$\tfrac{1}{2}$};
\node[vneg] (nm1_4) at (6.35,-3.15) {$-\tfrac{1}{4}$};
\node[vpos] (n3_1) at (11.41,-3.15) {$3$};
\node[vneg] (nm1_1) at (0.00,-4.20) {$-1$};
\node[vpos] (n3_2) at (2.19,-4.20) {$\tfrac{3}{2}$};
\node[vpos] (n3_4) at (6.35,-4.20) {$\tfrac{3}{4}$};
\node[vneg] (nm1_6) at (9.79,-4.20) {$-\tfrac{1}{6}$};
\node[vpos] (n4_1) at (13.02,-4.20) {$4$};
\node[vneg] (nm1_3) at (1.25,-5.25) {$-\tfrac{1}{3}$};
\node[vpos] (n5_2) at (3.12,-5.25) {$\tfrac{5}{2}$};
\node[vneg] (nm2_3) at (5.42,-5.25) {$-\tfrac{2}{3}$};
\node[vpos] (n7_4) at (7.29,-5.25) {$\tfrac{7}{4}$};
\node[vpos] (n5_6) at (9.79,-5.25) {$\tfrac{5}{6}$};
\node[vneg] (nm1_8) at (12.08,-5.25) {$-\tfrac{1}{8}$};
\node[vpos] (n5_1) at (13.96,-5.25) {$5$};
\node[vpos] (n2_3) at (1.25,-6.30) {$\tfrac{2}{3}$};
\node[vneg] (nm1_5) at (2.50,-6.30) {$-\tfrac{1}{5}$};
\node[vpos] (n7_2) at (3.75,-6.30) {$\tfrac{7}{2}$};
\node[vpos] (n1_3) at (5.42,-6.30) {$\tfrac{1}{3}$};
\node[vneg] (nm2_7) at (6.67,-6.30) {$-\tfrac{2}{7}$};
\node[vpos] (n11_4) at (7.92,-6.30) {$\tfrac{11}{4}$};
\node[vneg] (nm3_5) at (9.17,-6.30) {$-\tfrac{3}{5}$};
\node[vpos] (n11_6) at (10.42,-6.30) {$\tfrac{11}{6}$};
\node[vpos] (n7_8) at (12.08,-6.30) {$\tfrac{7}{8}$};
\node[vneg] (nm1_10) at (13.33,-6.30) {$-\tfrac{1}{10}$};
\node[vpos] (n6_1) at (14.58,-6.30) {$6$};
\node[vneg] (nm3_4) at (0.83,-7.35) {$-\tfrac{3}{4}$};
\node[vpos] (n5_3) at (1.67,-7.35) {$\tfrac{5}{3}$};
\node[vpos] (n4_5) at (2.50,-7.35) {$\tfrac{4}{5}$};
\node[vneg] (nm1_7) at (3.33,-7.35) {$-\tfrac{1}{7}$};
\node[vpos] (n9_2) at (4.17,-7.35) {$\tfrac{9}{2}$};
\node[vneg] (nm3_2) at (5.00,-7.35) {$-\tfrac{3}{2}$};
\node[vpos] (n4_3) at (5.83,-7.35) {$\tfrac{4}{3}$};
\node[vpos] (n5_7) at (6.67,-7.35) {$\tfrac{5}{7}$};
\node[vneg] (nm2_11) at (7.50,-7.35) {$-\tfrac{2}{11}$};
\node[vpos] (n15_4) at (8.33,-7.35) {$\tfrac{15}{4}$};
\node[vpos] (n2_5) at (9.17,-7.35) {$\tfrac{2}{5}$};
\node[vneg] (nm3_11) at (10.00,-7.35) {$-\tfrac{3}{11}$};
\node[vpos] (n17_6) at (10.83,-7.35) {$\tfrac{17}{6}$};
\node[vneg] (nm4_7) at (11.67,-7.35) {$-\tfrac{4}{7}$};
\node[vpos] (n15_8) at (12.50,-7.35) {$\tfrac{15}{8}$};
\node[vpos] (n9_10) at (13.33,-7.35) {$\tfrac{9}{10}$};
\node[vneg] (nm1_12) at (14.17,-7.35) {$-\tfrac{1}{12}$};
\node[vpos] (n7_1) at (15.00,-7.35) {$7$};
\node[rk] at (-0.6,0.63) {rank};
\node[rk] at (-0.6,0.00) {0};
\node[rk] at (-0.6,-1.05) {1};
\node[rk] at (-0.6,-2.10) {2};
\node[rk] at (-0.6,-3.15) {3};
\node[rk] at (-0.6,-4.20) {4};
\node[rk] at (-0.6,-5.25) {5};
\node[rk] at (-0.6,-6.30) {6};
\node[rk] at (-0.6,-7.35) {7};
\draw[ef] (n0_1) -- (n1_1);
\draw[eg] (n1_1) -- (nm1_2);
\draw[ef] (n1_1) -- (n2_1);
\draw[ef] (nm1_2) -- (n1_2);
\draw[eg] (n2_1) -- (nm1_4);
\draw[ef] (n2_1) -- (n3_1);
\draw[eg] (n1_2) -- (nm1_1);
\draw[ef] (n1_2) -- (n3_2);
\draw[ef] (nm1_4) -- (n3_4);
\draw[eg] (n3_1) -- (nm1_6);
\draw[ef] (n3_1) -- (n4_1);
\draw[eg] (n3_2) -- (nm1_3);
\draw[ef] (n3_2) -- (n5_2);
\draw[eg] (n3_4) -- (nm2_3);
\draw[ef] (n3_4) -- (n7_4);
\draw[ef] (nm1_6) -- (n5_6);
\draw[eg] (n4_1) -- (nm1_8);
\draw[ef] (n4_1) -- (n5_1);
\draw[ef] (nm1_3) -- (n2_3);
\draw[eg] (n5_2) -- (nm1_5);
\draw[ef] (n5_2) -- (n7_2);
\draw[ef] (nm2_3) -- (n1_3);
\draw[eg] (n7_4) -- (nm2_7);
\draw[ef] (n7_4) -- (n11_4);
\draw[eg] (n5_6) -- (nm3_5);
\draw[ef] (n5_6) -- (n11_6);
\draw[ef] (nm1_8) -- (n7_8);
\draw[eg] (n5_1) -- (nm1_10);
\draw[ef] (n5_1) -- (n6_1);
\draw[eg] (n2_3) -- (nm3_4);
\draw[ef] (n2_3) -- (n5_3);
\draw[ef] (nm1_5) -- (n4_5);
\draw[eg] (n7_2) -- (nm1_7);
\draw[ef] (n7_2) -- (n9_2);
\draw[eg] (n1_3) -- (nm3_2);
\draw[ef] (n1_3) -- (n4_3);
\draw[ef] (nm2_7) -- (n5_7);
\draw[eg] (n11_4) -- (nm2_11);
\draw[ef] (n11_4) -- (n15_4);
\draw[ef] (nm3_5) -- (n2_5);
\draw[eg] (n11_6) -- (nm3_11);
\draw[ef] (n11_6) -- (n17_6);
\draw[eg] (n7_8) -- (nm4_7);
\draw[ef] (n7_8) -- (n15_8);
\draw[ef] (nm1_10) -- (n9_10);
\draw[eg] (n6_1) -- (nm1_12);
\draw[ef] (n6_1) -- (n7_1);
\draw[eold] (nm3_2) -- (nm1_2);
\end{tikzpicture}
}

\medskip
{\small (b) $m=2$.}
\caption{Ranks $0$ to $7$ of the trees of $f(x)=x+1$ (\textcolor{fred}{red} edges) and
$g_m(x)=-1/(mx)$ (\textcolor{gblue}{blue} edges) for $m=1$ and $m=2$. Positive vertices are entered
by $f$ and negative ones by $g_m$ (Theorem~C$_q$). The vertices $x\le-1$ are leaves. A dashed arrow
shows $f$ applied to a leaf $x<-1$. The value $x+1$ is not new and its rank is lower by $2q-3$, which
is $3$ for $m=1$ and $5$ for $m=2$. The numbers on the left are ranks.}
\label{fig:trees}
\end{figure}

\subsection{The trees for every \texorpdfstring{$q$}{q}}

The proofs do not use $m$ at all. They use only the order $q$. With $X=\sqrt m\,x$ the maps become
$X\mapsto X+\sqrt m$ and $X\mapsto-1/X$, and for $m=1,2,3$ we have $\sqrt m=2\cos(\pi/q)$. So for
every $q\ge3$ we put
\[
\lambda=\lambda_q=2\cos\frac{\pi}{q},\qquad F(X)=X+\lambda,\qquad G(X)=-\frac1X ,
\]
and study the tree of $F$ and $G$ with root $0$, which we call the \emph{$\lambda_q$-tree}. For
$q=3,4,6$ it is the tree of $f$ and $g_m$ with $m=1,2,3$, scaled by $\sqrt m$
(Lemma~\ref{lem:coords}). For $q=5$ it lives in $\Q(\sqrt5)$, and for most larger $q$ in a number
field of higher degree. The maps $F$ and $G$ generate Hecke's group $H_q$ \cite{hecke1936}, and
$G\circ F$ has order $q$.

We write $d(X)$ for the rank of a point $X$ of the $\lambda_q$-tree, $a_q(n)$ for the number of points
of rank $n$, and $r_q(n)$ for the number of positive ones. We also put
\[
D_q(t)=1-t-t^3-t^5-\dots-t^{2q-3},\qquad
N_q(t)=1+t^2+t^4+\dots+t^{2q-4}-t^{2q-3}.
\]
Theorems A$_q$, B$_q$ and C$_q$ below hold for every $q\ge3$, and at $q=3$ they give the
corresponding parts of Theorems A, B(ii) and C of \cite{P137}.

\subsection{Results}

A string of integers $(c_1,\dots,c_k)$ with $k\ge1$, $c_1\ge0$ and $c_2,\dots,c_k\ge1$ is
\emph{reduced} if no $q-2$ consecutive digits among $c_2,\dots,c_k$ equal $1$, with one exception. If
$c_1=0$, the run of $1$'s that starts at $c_2$ may have length exactly $q-2$. Its \emph{word} is
$f^{c_1}gf^{c_2}g\cdots gf^{c_k}$, applied to $0$ from right to left, and its \emph{weight} is
$c_1+\dots+c_k+k-1$, the length of the word. For $q=3$ a reduced string has $c_2,\dots,c_k\ge2$, except
that $c_2=1$ is allowed when $c_1=0$. This is the rule of \cite[Theorem~B]{P137}.

\begin{thmB}[the rank formula, Theorem~\ref{thm:Bq}]
Let $q\ge3$. Every point of the $\lambda_q$-tree is reached by the word of exactly one reduced string,
and its rank is the weight of that string. This word is the only shortest word for the point. In
coordinates, the string $(c_1,\dots,c_k)$ gives
\[
X=c_1\lambda-\cfrac{1}{c_2\lambda-\cfrac{1}{\ddots-\cfrac{1}{c_k\lambda}}}\,,
\]
read as $-1/(\cdots)$ when $c_1=0$. For $m=2,3$ every rational $x$ has exactly one reduced string,
$x=c_1+g_m\bigl(c_2+g_m(\cdots+g_m(c_k))\bigr)$, and
\[
d(x)=c_1+c_2+\dots+c_k+k-1 .
\]
\end{thmB}

For example, take $m=2$ and $x=\tfrac35$. Then
\[
\tfrac35=1-\cfrac{1}{4-\cfrac{1}{1-\cfrac{1}{4-\cfrac11}}}\,,
\]
which is $1+g_2(2+g_2(1+g_2(2+g_2(1))))$. Its string is $(1,2,1,2,1)$, which is reduced for $q=4$
because no two consecutive digits after the first equal $1$. So $d(\tfrac35)=7+4=11$.

\begin{thmA}[the counts, Theorem~\ref{thm:Aq}]
Let $q\ge3$. Then
\[
\sum_{n\ge0}a_q(n)\,t^n=\frac{N_q(t)}{D_q(t)},
\]
and this fraction is in lowest terms. Hence
\[
a_q(n)=a_q(n-1)+a_q(n-3)+a_q(n-5)+\dots+a_q(n-2q+3)\qquad\text{for all }n\ge2q-2,
\]
and this fails at $n=2q-3$, where the left side minus the right side is $-1$. The positive points
have $\sum_n r_q(n)t^n=t(1+t^2+\dots+t^{2q-6})/D_q(t)$, and $a_q(n)=r_q(n)+r_q(n-1)$ for $n\ge1$.
\end{thmA}

The denominator depends only on $q$, and so does the whole count. For $m=2$ the positive counts
are $r(n)=\mathrm{A060961}(n-1)+\mathrm{A060961}(n-3)$, where A060961 counts compositions into
parts $1$, $3$ and $5$ \cite{oeisA060961}.

\begin{thmC}[the tree, Corollary~\ref{cor:Cq}]
Let $q\ge3$ and define the parent map $P(X)=X-\lambda$ for $X>0$ and $P(X)=-1/X$ for $X<0$. For
every point $X\neq0$ of the $\lambda_q$-tree, $P(X)$ is the only point of rank $d(X)-1$ from which
one move reaches $X$, and iterating $P$ reaches $0$ in exactly $d(X)$ steps. The last move of every
shortest word for $X$ is $F$ if $X>0$ and $G$ if $X<0$. So in Kagey's coloring, where a vertex is blue
when it is reached last by $G$, a vertex is blue if and only if it is negative. A
vertex $X>0$ has the two children $X+\lambda$ and $-1/X$, a vertex in $(-\lambda,0)$ has the one child
$X+\lambda$, and a vertex $X\le-\lambda$ is a leaf. For $X<-\lambda$ the value $X+\lambda$ appeared
exactly $2q-3$ ranks earlier.
\end{thmC}

In $x$-coordinates the leaves are the $x\le-1$, and for $x<-1$ the value $x+1$ appeared five ranks
earlier when $m=2$ and nine ranks earlier when $m=3$. This jump of $2q-3$ ranks is where the lag
$2q-3$ of the recurrence comes from.

\begin{thmD}[reachability, Theorem~\ref{thm:reach}]
For $m=1,2,3$ every rational number is in the tree of $f(x)=x+1$ and $g_m(x)=-1/(mx)$.
\end{thmD}

For $m=1$ this is \cite[Theorem~3.6]{P137}. The group generated by $f$ and $g_m$ is known to move $0$
to every rational for $m=2,3$ (Remark~\ref{rem:group}). Theorem~D is stronger, because the tree uses
only the forward moves $f$ and $g_m$ and never $f^{-1}$.

\begin{thmE}[the free case, Theorem~\ref{thm:free}]
Let $m\ge4$. The tree of $f(x)=x+1$ and $g_m(x)=-1/(mx)$ has $a(n)=F_{n+1}$ points of rank $n$,
where $F_n$ is the Fibonacci sequence with $F_1=F_2=1$. Every finite string $(c_1,\dots,c_k)$ with
$c_1\ge0$ and $c_2,\dots,c_k\ge1$ (a free string) gives a point, distinct strings give distinct
points, and the rank of the point is the weight of the string. A vertex is blue if and only if it is negative, and no
rational in $(-\infty,-2/m]$ is reached. In particular $-1$ and $\tfrac12$ are never reached. This
covers the tree of $x\mapsto x+k$ and $x\mapsto-1/x$ for every $k\ge2$.
\end{thmE}

The growth rate of $a_q(n)$ is the largest root $\psi_q$ of
$\psi^{2q-3}=\psi^{2q-4}+\psi^{2q-6}+\dots+1$. These numbers increase with $q$ and tend to the
golden ratio (Proposition~\ref{prop:growth}).

\subsection{How the proofs go}

The proofs follow \cite{P137} step by step. Section~\ref{sec:parent} proves that the parent map
traces a shortest path, which gives Theorem~C$_q$. The key step is the analogue of the identity
$d(x-1)=d(x)+3$ for $x<0$, which becomes $d(X-\lambda)=d(X)+2q-3$. Its proof follows a path of $2q-2$
steps through the orbit of $\infty$ under $F\circ G$, which has length $q$. Section~\ref{sec:reach}
proves Theorem~D with a height that drops when we shift $x$ to the nearest integer and apply $g_m$. Section~\ref{sec:words} proves
Theorem~B$_q$ and explains the run rule. Section~\ref{sec:count} counts reduced strings and proves
Theorem~A$_q$, and Section~\ref{sec:free} treats $m\ge4$. Section~\ref{sec:verify} lists the
computer checks, and Section~\ref{sec:final} discusses related work and open questions. None of the
proofs depends on the computations or on facts about Hecke groups.
